\section{Open Questions and Recommended Next Steps}

Five questions this replication surfaces but does not answer. Each is a concrete, tractable follow-up that would sharpen or extend the paper's contribution.

\subsection*{Q1. Discriminating experiment on a genuinely reducible group}
Does the França--Hashagen framework produce demonstrably better fidelity extraction than standard Clifford RB when applied to a hardware-native gate group whose representation is genuinely reducible (e.g., a 2-qubit iSWAP-family group, or a qutrit XX+YY gate group), where a single-exponential fit provably fails?

\emph{Basis.} The paper's key theoretical novelty --- a multi-term exponential fit whose degree equals the number of inequivalent irreps of $G$ --- was not exercised by our replication. We used single-exponential fits that would also succeed for exact 2-designs. Section 6.2 of the paper mentions this generalisation but does not include a discriminating experiment.

\emph{Next steps.} Pick a specific reducible finite group $G$ (e.g., $\langle \mathrm{iSWAP}, S, X\rangle$ on 2 qubits, or a Cirq/Qiskit-supported hardware-native family). Enumerate irreps via GAP (free) or hand character-table analysis. Run RB fitted to (i) a naive single-exponential model and (ii) the paper's $k$-term irrep-count model. Compare fit residuals and recovered $F$ against analytic ground truth. A $\ge 2\times$ residual reduction from (i) to (ii) would be a discriminating positive result. Wall time: $\sim 4$--$8$\,h on a single CPU node.

\subsection*{Q2. Empirical tightness of concentration bounds}
Are the paper's concentration bounds for the RB survival-probability estimator (finite-$M$ sample-complexity regime) empirically tight, or does the true finite-sample tail probability decay faster/slower than the derived bound?

\emph{Basis.} We reported only point-estimator errors, not tail-probability sweeps. The paper's theorem statements give upper bounds on $\Pr[|\hat P_m - P_m| > \varepsilon]$ but do not report empirical tail-probability curves. Whether the constants are tight or loose changes practical resource-estimates for lab-scale RB campaigns.

\emph{Next steps.} Fix $(d, m, p) = (8, 20, 0.9)$. Sweep $M \in \{10, 30, 100, 300, 1000, 3000\}$ with $500$ independent replicates per $M$ (parallelizable across m1 / spark mesh). Compute empirical $\Pr[|\Delta| > \varepsilon]$ for $\varepsilon \in \{10^{-4}, 10^{-3}, 10^{-2}\}$ and compare to the paper's Hoeffding-style bound. Fit a data-driven exponent; if empirical scaling is faster, the bound is loose and can be tightened.

\subsection*{Q3. Robustness under group-structure-breaking noise}
How does the finite-group RB estimator behave under noise that explicitly breaks the group structure --- e.g., non-$\delta$-covariant noise, or noise correlated with the specific group element being applied?

\emph{Basis.} Both the paper and our replication assumed $\delta$-covariant noise (randomly rotated depolarising channel), which respects the RB decoupling. The framework's robustness under structured non-covariant noise is not established. Real-hardware crosstalk and gate-dependent noise systematically violate this assumption.

\emph{Next steps.} Construct a controlled gate-dependent noise family $T_g(\rho) = p\rho + (1-p) U_g \rho U_g^\dagger$ where $U_g$ is a bounded perturbation of the ideal $g$, with perturbation strength $\varepsilon \in [0, 0.1]$. Run RB on $\mathrm{MU}(4,8)$ (fits inside our existing harness --- modify only the noise-application step) and measure $|\hat F - F_{\text{true}}|$ vs $\varepsilon$. Compare to the $O(\varepsilon)$ bias predicted by first-order perturbation analysis. Reuse existing \texttt{monomial\_rb.py}, $\sim 2$\,h wall on m1.

\subsection*{Q4. ML-boosted estimator variants}
Can an ML-boosted estimator (learned-basis regression, or a small neural point estimator) recover fidelity from shorter/fewer sequences than the paper's least-squares fit, at matched estimation error?

\emph{Basis.} The paper's estimator is a two-parameter exponential fit via least-squares. Modern learned estimators (simulation-based inference, neural point estimators) have achieved 2--5$\times$ sample-complexity reduction in adjacent noise-characterization work. Whether this holds for finite-group RB --- where the fit function is analytically well-specified --- is open.

\emph{Next steps.} Generate $10^4$--$10^5$ synthetic RB trajectories from \texttt{monomial\_rb.py} with $p$ and channel drawn from a wide prior; train a small MLP (or GBM) to map $\{P_m\}_{m \in \text{grid}} \to \hat F$. Benchmark against the exponential-fit estimator at held-out $(p, \sigma)$. Metric: matched MSE at what fraction of the $M$ budget? Free stack: scikit-learn + local PyTorch, $\sim 4$\,h on m1, no GPU needed.

\subsection*{Q5. Extension to continuous gate sets}
Does the finite-group RB framework extend meaningfully to continuous gate sets (parameterised rotations) via a discretisation / $\varepsilon$-net argument, or is the discreteness of $G$ essential?

\emph{Basis.} The paper requires $G$ finite. Real hardware primitives (arbitrary-angle rotations, cross-resonance gates) are continuous. Naive $\varepsilon$-net discretisation of SU$(d)$ produces exponentially many irreps and blows up fit complexity. Whether a smarter discretisation (e.g., a $T$-design cover) preserves tractability is not addressed.

\emph{Next steps.} Literature review (arxiv-sanity, semantic-scholar free tier): search 'continuous randomized benchmarking', 'gate-set RB', 'approximate design RB' (2019--2025). Then minimal numerical probe: $\{R_z(2\pi k / N) : k = 0, \dots, N-1\}$ for $N \in \{8, 16, 32, 64\}$, treat as $\mathbb{Z}_N$, run MU-style RB, study convergence of $\hat F$ as $N \to \infty$. If $\hat F$ converges to the continuous-family fidelity, discretisation works. $\sim 1$ day of Ollie time using existing infra.
